\subsection{הבעיה}

בתחרויות \textenglish{FRC} משתתפות בכל מקצה שתי בריתות — כחולה
ואדומה — כאשר בכל ברית שלושה רובוטים של קבוצות שונות. במהלך המשחק כל
רובוט נע במהירות, נבלע לעיתים מאחורי רובוטים אחרים, ומשתנה בזווית
הצילום שלו ביחס למצלמה. כדי לבחור שותפות טובות למקצי הפלייאוף ולבצע החלטות אסטרטגיות אחרות, קבוצות
רבות מבצעות \textenglish{Scouting}: תלמידים צופים במקצים ומזינים ידנית
נתונים על ביצועי הרובוטים.

לשיטה הידנית יש כמה חסרונות בולטים: היא דורשת כוח אדם רב, היא לרוב לא מאפשרת מעקב מלא אחרי מיקומי הרובוטים, והיא רגישה לטעויות אנוש. לכן
נוצר צורך במערכת אוטומטית שתוכל לקבל סרטון וידאו מהיציע, לזהות את כל
הרובוטים במקצה, לעקוב אחרי התנועה שלהם לאורך זמן, ולשייך כל מסלול
לקבוצה הנכונה.

\subsection{הפתרון המוצע}

המערכת שפיתחתי מקבלת כקלט סרטון של מקצה
ומספר מקצה ידוע מראש. מתוך מספר המקצה המערכת יודעת מהן שש הקבוצות
האפשריות שישתתפו בו, ולכן אפשר לצמצם את מרחב החיפוש כבר בתחילת
התהליך. לאחר מכן המערכת פועלת בארבעה שלבים עיקריים:

\begin{itemize}
\item זיהוי כל הרובוטים בפריים באמצעות מודל \textenglish{YOLO}.
\item מעקב אחר כל זיהוי לאורך זמן באמצעות \textenglish{ByteTrack}, כך
  שלכל רובוט נשמר \textenglish{Track ID} יציב ככל האפשר.
\item חיתוך תמונת הרובוט מתוך הפריים והרצת מודל \textenglish{YOLO}
  נוסף לזיהוי הספרות שעל הבאמפר.
\item התאמת כל עקבה למספר קבוצה מתוך שש הקבוצות במקצה בעזרת מנגנון
  ניקוד מצטבר, המשלב את צבע הברית, מרחק
  \textenglish{Levenshtein} ממספרי הקבוצות האפשריים, וסף ביטחון לפני
  קבלת החלטה סופית.
\end{itemize}

פלט המערכת הוא מסלול התנועה המוחלק של כל רובוט לאורך כל המקצה. בצורה זו ניתן לעבור משלב של צפייה ידנית בלבד לשלב של ניתוח
אוטומטי למחצה או אוטומטי מלא של התנהגות הרובוטים לאורך המשחק.

\subsection{פתרונות דומים והייחוד של הפרויקט}

הגישה הנפוצה ביותר כיום בסביבת \textenglish{FRC} היא עדיין סקאוטינג
ידני, שבו תלמידים עוקבים אחרי רובוטים ומזינים נתונים בזמן אמת. קיימות
גם מערכות ראייה ממוחשבת כלליות לזיהוי ומעקב אחר אובייקטים, אך הן בדרך
כלל אינן פותרות לבדן את שאלת הזהות של כל רובוט כאשר כמה רובוטים דומים
זה לזה ומופיעים באותו מגרש.

הייחוד של הפרויקט שלי הוא בשילוב בין כמה מקורות מידע חלשים יחסית,
שביחד יוצרים החלטה חזקה: זיהוי רובוטים, מעקב תנועתי, זיהוי ספרות על
הבאמפר, וצמצום המועמדים האפשריים לפי לוח המשחקים של המקצה. במקום לנסות
לזהות את מספר הקבוצה במדויק בכל פריים בנפרד, המערכת צוברת ראיות לאורך
זמן ומקבלת החלטה רק כאשר הביטחון המצטבר מספיק גבוה.

\subsection{סרטון הדגמה}

יש להוסיף כאן קישור לסרטון הדגמה קצר של המערכת בפעולה
(זיהוי רובוטים, מעקב ושיוך קבוצות): \textbf{[להוסיף קישור]}.

\subsection{סיכום תוצאות}

בבדיקות שביצעתי התקבלו התוצאות העיקריות הבאות:
\begin{itemize}
\item מודל זיהוי הרובוטים השיג \textenglish{mAP@50} של כ־\textbf{0.95}.
\item מודל זיהוי הספרות הגיע לדיוק של כ־\textbf{0.70} ברמת הספרה
  הבודדת.
\item לאחר שילוב זיהוי הספרות עם צבע הברית, לוח המשחקים ואלגוריתם
  \textenglish{Levenshtein}, הדיוק הסופי בשיוך עקבה לקבוצה עלה
  ל־\textbf{[להשלים]}.
\item המערכת מצליחה להפיק עבור כל קבוצה מסלול תנועה רציף, שניתן להציג
  לאחר העיבוד לצורך ניתוח המשחק.
\end{itemize}

שתי תוצאות אלו מדגישות את הרעיון המרכזי של הפרויקט: גם כאשר זיהוי
הספרות לבדו עדיין אינו מושלם, שילוב מושכל של כמה שכבות מידע מאפשר
לשפר משמעותית את ההחלטה הסופית ולקבל מערכת שימושית לסקאוטינג.
